% \documentclass[compacto,10pt]{aleph-notas}
\documentclass[a4,10pt,comentarios]{aleph-notas}

% -- 额外宏包
\usepackage{enumitem}
\usepackage{aleph-comandos}

% -- 文档信息
\institucion{武汉大学}
\carrera{机器人工程}
\asignatura{材料力学}
\tema{笔记模板}
\autor{Lai Wei}
\date{\today}

%% --> 校徽
\logouno[5cm]{Logos/whulogo.png}
\definecolor{colordef}{cmyk}{0.81,0.62,0.00,0.22}

\begin{document}

\encabezado

%%%%%%%%%%%%%%%%%%%%%%%%%%%%%%%%%%%%%
\section{平面直角坐标系}
%%%%%%%%%%%%%%%%%%%%%%%%%%%%%%%%%%%%%

%%%%%%%%%%%%%%%%%%%%%%%%%%%%%%%%%%%%%
%%%%%%%%%%%%%%%%%%%%%%%%%%%%%%%%%%%%%
\begin{defi}
    平面直角坐标系是集合
    \[
        \R^2 = \{(x,y): x,y\in\R\}
    \]
    的图形表示。
\end{defi}

\begin{teo}[三角形面积]
    若三角形的顶点为 $(x_1,y_1),(x_2,y_2),(x_3,y_3)\in\R^2$，则该三角形的面积为
    \[
        \frac{1}{2}
        \begin{vmatrix}
            x_1 & y_1 & 1 \\
            x_2 & y_2 & 1 \\
            x_3 & y_3 & 1 \\
        \end{vmatrix}
    \]
    的绝对值。
\end{teo}

\begin{defi}[线段]
    对于 $(x_1,y_1),(x_2,y_2)\in\R^2$，令
    \[
        A=(x_1,y_1)
        \texty
        B=(x_2,y_2),
    \]
    则形如
    \[
        A + t(B-A) = (x_1 + t(x_2-x_1),\ y_1 + t(y_2-y_1)),
    \]
    其中 $t\in[0,1]$ 的点，构成连接 $A$ 与 $B$ 的线段。
\end{defi}



\begin{comentarios}
    此处可通过以下活动引入 GeoGebra 动画。

    \begin{ejer}
        定义两个点，并嵌入一个从一点移动到另一点的文字。
    \end{ejer}
    只需利用之前制作的活动，使用\emph{文本}工具在点 $C$ 上添加文字即可。
\end{comentarios}


\begin{advertencia}
    若在上述定义中取 $t=1/2$，则点
    \[
        A + \tfrac{1}{2}(B-A) = \frac{1}{2}(A+B)
        =\left(\frac{x_1+x_2}{2},\frac{y_1+y_2}{2} \right)
    \]
    为 $A$ 与 $B$ 的中点。
\end{advertencia}




\end{document}
